\subsubsection{Independent holomorphic shielding envelopes}
\label{subsec:holomorphic-shielding}

The signed budget $4\int_0^\sigma \mathcal D=G(0)-G(\sigma)$
is an endpoint identity, not an independent upper estimate.
A direct holomorphic contour bound supplies such an estimate
on a genuine, albeit initially short, theta-flow interval.

\begin{theorem}[Local Rouch\'e-to-shielding transfer]
\label{thm:rouche-shielding-independent}
Fix $t_*>\Lambda$, a time at which the de Bruijn--Newman
theta transform
$H_t(z)=\int_0^\infty \Phi(u)e^{tu^2}\cos(zu)\,du$
has simple real zeros. Let $a<b$ be consecutive zeros of
$H_{t_*}$, with $g_*=b-a$. Choose $0<r<g_*/2$
so that the two disks $D_a,D_b$ of radius $r$, centred
at $a,b$, are disjoint and each contains one zero.
Choose a conjugation-invariant Jordan domain $\Omega$
containing both disks and the real segment $[a-r,b+r]$,
with exactly two zeros in its interior and no zeros
on its boundary. Put
\[
Y=\max_{\partial D_a\cup\partial D_b\cup\partial\Omega}
|\Im z|,\quad
m=\min_{\partial D_a\cup\partial D_b\cup\partial\Omega}
|H_{t_*}(z)|>0,
\]
and
\[
\mathcal M_Y=\int_0^\infty
\Phi(u)e^{t_*u^2}u^2e^{Yu}\,du<\infty.
\]
If $0<\sigma\le t_*$ and $\sigma\mathcal M_Y<m$,
then for all $0\le\tau\le\sigma$ the function
$H_{t_*-\tau}$ has one simple real zero in each
disk, with no intervening zero. Their separation
obeys $g(\tau)\ge g_*-2r>0$.

On any such interval on which the classical real-simple
zero-motion law $G'=-4\mathcal D$ holds,
\begin{equation}
\boxed{
\int_0^\tau\mathcal D(v)\,dv
\le\mathcal B_r=r(g_*-r),\qquad
4\mathcal B_r=g_*^2-\eta_r,\quad
\eta_r=(g_*-2r)^2>0.}
\label{eq:rouche-independent-shielding}
\end{equation}
The contour hypothesis is independent of
the unknown moving gap.
\end{theorem}
\begin{proof}
For $u\ge0$, $\Phi(u)>0$,
$1-e^{-\tau u^2}\le\tau u^2$,
and $|\cos(zu)|\le e^{|\Im z|u}$. Hence
on every selected contour,
\[
|H_{t_*-\tau}(z)-H_{t_*}(z)|
\le\tau\mathcal M_Y\le\sigma\mathcal M_Y<m.
\]
Rouch\'e preserves one zero in each disk
and two in $\Omega$, counted with multiplicity.
Conjugation symmetry makes each unique disk zero
real and simple, while the enclosing contour excludes
another zero between them. This proves the separation
bound. Only now invoke the signed gap ODE:
\[
4\int_0^\tau\mathcal D
=G(0)-G(\tau)
\le g_*^2-(g_*-2r)^2=4r(g_*-r).
\]
The last equality and the strictly positive margin
follow by algebra.
\end{proof}

\begin{corollary}[Pair-dependent nonzero theta interval]
\label{cor:local-rouche-shield}
For every fixed consecutive pair of zeros of
$H_{1/2}$, an $r>0$ and $\sigma_i>0$
satisfying the preceding theorem exist.
The same holds at $t_*=1/4$, since the
published unconditional bound $\Lambda\le0.22$
implies $1/4>\Lambda$.
The explicit contour threshold $m/\mathcal M_Y$ obtained here need not be uniform in height. Independently, the published bound $\Lambda\le0.22$ already guarantees an all-height real-simple slab for $t>0.22$; the new local contour estimate is useful as an independent, computable signed-defect certificate.
\end{corollary}
\begin{proof}
The initial zeros are isolated and simple.
Choose disjoint disks and an enclosing contour
with strictly positive boundary minima.
The theta moment $\mathcal M_Y$ is finite by
superexponential kernel decay, so one may take
$0<\sigma_i<\min(t_*,m/\mathcal M_Y)$.
\end{proof}

\paragraph{An exact adversarial control.}
For the quartic backward heat flow
$p_s=(z^2-1)(z^2-4)+s(12z^2-10)+12s^2$,
choose four disks of radius $r=1/4$ around
$-2,-1,1,2$. On each boundary
\[
|p_0|\ge\tfrac14\tfrac34\tfrac74\tfrac{11}4
=\tfrac{231}{256},
\quad |p_s-p_0|\le\tfrac{283}{4}s+12s^2.
\]
For $s\le1/80$ the latter is at most
$709/800$, and the exact margin is
$231/256-709/800=103/6400>0$.
Rouch\'e therefore keeps one real root per disk.
For the protected positive pair this yields
\[
g(0)=1,\qquad\mathcal B=3/16,\qquad
\eta=1/4,\qquad
4\int_0^{1/80}\mathcal D\le3/4.
\]
Nevertheless the quartic has a later collision at
$(10-\sqrt{46})/24>1/80$. Local shielding
certification cannot be extrapolated to the
full backward heat interval.

\paragraph{Proof boundary.}
For the actual theta family the theorem provides a
noncircular local margin for every fixed pair.
It does not construct boundary minima uniform in
pair index or throughout $t\downarrow0$, and does
not rule out zeros arriving from unbounded height.
Those are still the RH-level obligations.
The rational transfer and the quartic contour
are checked in \path{XiKernel/RoucheShieldingTransfer.lean}
and \path{certificates/xi_rouche_shielding_local_exact.py};
the full theta Fourier-integral and Rouch\'e argument
are manuscript proofs, not formalized end-to-end in Lean.
